
\documentclass[border=4pt]{standalone}
\usepackage{polyglossia}
\setdefaultlanguage{russian}
\setotherlanguage{english}
\usepackage{fontspec}
\IfFontExistsTF{CMU Serif}{\setmainfont{CMU Serif}}{\setmainfont{Liberation Serif}}
\IfFontExistsTF{CMU Sans Serif}{\setsansfont{CMU Sans Serif}}{\setsansfont{Liberation Sans}}
\IfFontExistsTF{CMU Typewriter Text}{\setmonofont{CMU Typewriter Text}}{\setmonofont{DejaVu Sans Mono}}
\usepackage{amsmath,amssymb,mathtools,bm}
\usepackage{xcolor}
\definecolor{accent}{HTML}{1F4E79}
\definecolor{accent2}{HTML}{B85450}
\definecolor{shade}{HTML}{F4F1EA}
\colorlet{prosvBlue}{accent}
\colorlet{prosvRed}{accent2}
\colorlet{prosvLight}{accent!12}
\colorlet{prosvCream}{shade}
\definecolor{prosvGold}{HTML}{E8B647}
\definecolor{prosvGreen}{HTML}{6A9F58}
\colorlet{prosvGray}{black!55}
\usepackage{tikz}
\usetikzlibrary{calc, arrows.meta, decorations.pathreplacing,
                positioning, shapes.geometric, shapes.misc, fit,
                backgrounds}
\usepackage{pgfplots}
\pgfplotsset{compat=1.18}
\newcommand{\R}{\mathbb{R}}
\newcommand{\N}{\mathbb{N}}
\newcommand{\Z}{\mathbb{Z}}
\newcommand{\eps}{\varepsilon}
\newcommand{\norm}[1]{\left\|#1\right\|}
\newcommand{\abs}[1]{\left|#1\right|}
\newcommand{\NOD}{\gcd}
\newcommand{\Eul}{\varphi}
\newcommand{\Zn}[1]{\mathbb{Z}_{#1}}
\newcommand{\modn}[1]{\,(\mathrm{mod}\,#1)}
\newcommand{\term}[1]{\textcolor{accent2}{\textbf{#1}}}
\newcommand{\engterm}[1]{\textit{#1}}
\DeclareMathOperator*{\argmin}{arg\,min}
\DeclareMathOperator*{\argmax}{arg\,max}
\begin{document}

\begin{tikzpicture}[
  font=\small,
  every node/.style={align=center},
  alice/.style={fill=prosvLight, draw=prosvBlue, rounded corners=2pt, minimum width=4.5cm, minimum height=0.7cm},
  bob/.style={fill=prosvCream, draw=prosvGold!80!brown, rounded corners=2pt, minimum width=4.5cm, minimum height=0.7cm}
]
  
  \node[font=\bfseries\color{prosvBlue}] at (-3.5, 4) {Алиса};
  \node[font=\bfseries\color{prosvRed}] at (3.5, 4) {Боб};
  \node[font=\bfseries\color{prosvGreen}] at (0, 4) {Открытый канал};
  
  
  \node[alice] (A1) at (-3.5, 3) {Выбирает секрет $a$};
  \node[alice] (A2) at (-3.5, 2) {Считает $A = g^a \bmod p$};
  \node[alice] (A3) at (-3.5, 0.5) {Получает $B$};
  \node[alice] (A4) at (-3.5, -0.5) {$K = B^a \bmod p$};
  
  
  \node[bob] (B1) at (3.5, 3) {Выбирает секрет $b$};
  \node[bob] (B2) at (3.5, 2) {Считает $B = g^b \bmod p$};
  \node[bob] (B3) at (3.5, 0.5) {Получает $A$};
  \node[bob] (B4) at (3.5, -0.5) {$K = A^b \bmod p$};
  
  
  \draw[->, thick, prosvBlue] (A2) -- node[above, font=\scriptsize]{$A$} (B3);
  \draw[->, thick, prosvRed]  (B2) -- node[below, font=\scriptsize]{$B$} (A3);
  
  
  \node[fill=prosvLight!30!white, draw=prosvGreen, very thick, rounded corners=3pt, minimum width=7cm, minimum height=0.8cm] at (0,-1.8) {Общий ключ $K = g^{ab} \bmod p$};
\end{tikzpicture}
\end{document}
